\begin{enumerate}
    \item Einleitung
    \begin{itemize}
        \item Problemstellung und Motivation
        \item Zielsetzung der Arbeit
        \item Aufbau der Arbeit
    \end{itemize}

    \item Grundlagen und Stand der Forschung
    \begin{itemize}
        \item Begriffliche Grundlagen
        \item Bestehende Ansätze und Abgrenzung
    \end{itemize}

    \item Konzeption und Umsetzung des Lösungsansatzes
    \begin{itemize}
        \item LLM-basierter Ansatz
        \item Regelbasierter Referenzansatz
    \end{itemize}

    \item Evaluierung der Ergebnisse
    \begin{itemize}
        \item Evaluationsmethodik und Datensatz
        \item Vergleich der Ansätze
    \end{itemize}

    \item Zusammenfassung und Ausblick
\end{enumerate}

\subsection*{Vorläufiger Zeit- und Aufgabenplan}

Die Bearbeitungszeit der Bachelorarbeit erstreckt sich über den Zeitraum von Mitte Mai bis Mitte August und umfasst insgesamt ca. 12 Wochen.

\begin{tabular}{p{3cm} p{10cm}}
\toprule
\textbf{Zeitraum} & \textbf{Aufgaben} \\
\midrule

Woche 1--2 &
Einarbeitung in das Themengebiet sowie Durchführung einer systematischen Literaturrecherche; Präzisierung der Anforderungen und Analyse des Datenbankschemas \\

Woche 3--4 &
Analyse der Ausgangsdaten und Definition der relevanten Datenfelder; Konzeption des LLM-basierten Ansatzes sowie des regelbasierten Referenzansatzes \\

Woche 5--7 &
Implementierung des LLM-basierten Mapping-Ansatzes; Entwicklung und Implementierung des regelbasierten Vergleichsansatzes \\

Woche 8--9 &
Erstellung des Evaluationsdatensatzes (Ground Truth) sowie Definition und Implementierung der Evaluationsmetriken \\

Woche 10--11 &
Durchführung der Evaluation; Analyse und Vergleich der Ergebnisse hinsichtlich Qualität, Schema-Konformität und Effizienz \\

Woche 12 &
Zusammenführung der Ergebnisse sowie Ausarbeitung der Kapitel zur Evaluierung und Diskussion \\

Woche 13 &
Finalisierung der Arbeit, Korrekturphase sowie Vorbereitung der Abgabe \\

\bottomrule
\end{tabular}